\documentclass[12pt,a4paper]{article}

% ===== PACKAGE IMPORTS =====
\usepackage[utf8]{inputenc}
\usepackage[english]{babel}
\usepackage[a4paper,left=17mm,top=22mm,right=17mm,bottom=15mm]{geometry}

% Math packages
\usepackage{amsmath,amsfonts,amssymb}
\usepackage{mathtools}
\usepackage{bm}
\usepackage[T1]{fontenc}
\usepackage{amsthm}

% Graphics and floats
\usepackage{graphicx}
\usepackage{float}
\usepackage{multirow}
\usepackage{booktabs}

\floatstyle{plain}
\newfloat{algorithm}{ht}{loa}
\floatname{algorithm}{Algorithm}

% Links and references
\usepackage[colorlinks=true, citecolor=blue, linkcolor=black, urlcolor=blue]{hyperref}
\usepackage{url}

% Code listings
\usepackage{listings}
\usepackage{xcolor}

\lstdefinelanguage{Python}{
  keywords={def,class,if,elif,else,for,while,break,continue,
            return,import,from,as,with,try,except,finally,raise,
            lambda,pass,yield,in,is,not,and,or,global,nonlocal,
            assert,True,False,None},
  keywordstyle=\color{blue}\bfseries,
  sensitive=true,
  comment=[l]{\#},
  commentstyle=\color{gray}\itshape,
  string=[b]",
  morestring=[b]',
  stringstyle=\color{red},
  basicstyle=\ttfamily\footnotesize,
  breaklines=true,
  tabsize=2,
  showstringspaces=false
}
\lstset{
  language=Python,
  frame=single,
  numbers=left,
  numberstyle=\tiny\color{gray},
  xleftmargin=2em,
  framexleftmargin=1.5em,
  basicstyle=\ttfamily\small,
  backgroundcolor=\color{gray!10},
  rulecolor=\color{black!30},
  captionpos=b,
  breaklines=true,
  tabsize=2,
  showstringspaces=false
}

% ===== THEOREM STYLES =====
\theoremstyle{plain}
\newtheorem{theorem}{Theorem}[section]
\newtheorem{proposition}[theorem]{Proposition}
\newtheorem{lemma}[theorem]{Lemma}
\newtheorem{corollary}[theorem]{Corollary}

\theoremstyle{definition}
\newtheorem{definition}[theorem]{Definition}
\newtheorem{assumption}[theorem]{Assumption}

\theoremstyle{remark}
\newtheorem{remark}[theorem]{Remark}

% ===== MACROS =====
\newcommand{\R}{\mathbb{R}}
\newcommand{\E}{\mathbb{E}}
\newcommand{\Tr}{\mathrm{Tr}}
\newcommand{\diag}{\mathrm{diag}}
\newcommand{\supp}{\mathrm{supp}}
\newcommand{\1}{\mathbf{1}}
\newcommand{\VDT}{\textsc{VDT}}

\title{Vibrational Deduction Transformers:\\
Laplacian Wave-Based Recurrent Transformer}

\author{Lorenzo Moriondo\\
\small Independent Researcher --- tuned.org.uk\\
\small ORCID: 0000-0002-8804-2963}

\date{2 June 2026}

\begin{document}

\maketitle

\begin{abstract}
We develop a self-contained theory of \emph{vibrational learning} on
feature space and use it to construct the \emph{Vibrational Deduction
Transformer} (\VDT{}), a forward-only recurrent architecture whose
intermediate states are constrained by the spectral geometry of a
feature-space graph Laplacian.
The foundations are classical and mechanical.
The transpose $A^\top$ of a data matrix transports residuals from sample
space to feature space; equipping that feature space with a graph
Laplacian $L_f$ built from $A^\top$ as a stiffness operator and a
positive diagonal mass matrix $M$ as inertia turns it into a
mass--spring system in the sense of Rayleigh.
The pair $(L_f, M)$ defines a generalised eigenproblem $L_f u = \lambda M u$
whose eigenvalues are squared natural frequencies and whose
eigenvectors are mode shapes, and the generalised Rayleigh quotient
$R_M(z) = z^\top L_f z / z^\top M z$ measures the spectral energy of any
direction.
From this footing we recover Laplacian-preconditioned gradient descent
with a provable linear-convergence guarantee, taumode frequency
regularisation that couples learning to the energy-dispersion structure
of a corpus, and a second-order damped-oscillator reading of gradient
descent itself.
We then lift these mechanics from a single parameter vector to the
hidden states of a recurrent transformer.
Each forward pass alternates an expressive transformer mixing block with
a structured \emph{vibrational state update} --- a discrete damped wave
operator $\Phi_L$ on the feature Laplacian --- so that every recurrent
step is forced through a physically interpretable spectral manifold,
in direct analogy to the lattice-projection step of the Lattice
Deduction Transformer~\cite{davis2026ldt}.
We further equip the state with a \emph{signed density matrix}
$\varrho_t = \varrho_t^{+} - \varrho_t^{-}$ (each component positive
semidefinite) in the Laplacian eigenbasis, capturing positive and
negative evidential mass, mode--mode coherence, and interferometric
cancellation.
We give the architecture, its spectral (modal) decoupling, the training
objectives, and an evaluation protocol mirroring the Lattice Deduction
Transformer benchmark exactly, with interpretable modal energy spectra
as a diagnostic.
\end{abstract}

%-------------------------------------------------------------------
\section{Introduction}
\label{sec:intro}
%-------------------------------------------------------------------

Reasoning in large neural networks has become largely synonymous with
the \emph{chain-of-thought} paradigm~\cite{wei2022cot} and its
test-time-compute extensions~\cite{snell2024scaling,brown2024large},
in which a model produces explicit intermediate tokens before
committing to a final answer.
The appeal is transparency: each reasoning step is a readable string,
and generating more tokens at inference time can improve
accuracy~\cite{muennighoff2025scaling}.
Yet the paradigm couples \emph{computation} to \emph{communication}:
every intermediate thought must be expressed as a sequence of
vocabulary items, and the cost of additional reasoning is proportional
to the number of output tokens rather than to the intrinsic difficulty
of the computation being performed.

An alternative --- closer in spirit to classical automated
reasoning --- is to treat inference as \emph{iterative refinement of a
structured latent state}, analogous to discrete diffusion over a
partial-information
representation~\cite{austin2021d3pm,davis2026ldt}.
Here, no intermediate tokens are emitted; instead the network
repeatedly updates a hidden state constrained to lie in a domain with
algebraic or geometric semantics.
Each update is a provably structured step --- a partial deduction in a
logical lattice~\cite{davis2026ldt}, a diffusion step on a discrete
space~\cite{austin2021d3pm} --- and the number of iterations trades
off compute against solution quality without touching the output
vocabulary.

This paper proposes a third, \emph{physically grounded} instance of
the same paradigm.
We equip the latent state with the spectral geometry of a
\emph{feature-space graph Laplacian} $L_f$, built from the transposed
data matrix $A^\top$, and constrain each recurrent step to propagate
along the eigenmodes of $(L_f, M)$ via a discrete damped wave
equation.
The Laplacian provides a \emph{discrete geometry} --- a spectral
lattice of natural frequencies --- and the vibrational operator
provides the \emph{forward approximation mechanism}, in direct analogy
to the role played by lattice projection in the Lattice Deduction
Transformer~\cite{davis2026ldt}.
Where chain-of-thought externalises reasoning into tokens, and where
lattice-deduction internalises it into logical candidate sets, the
\emph{Vibrational Deduction Transformer} (\VDT{}) internalises it
into the \emph{modal dynamics of feature space}: a physical wave
settles toward an answer the way a struck string settles toward
silence, each recurrent step dissipating energy in spectrally
structured fashion.

Two ideas motivate this paper, one about \emph{learning dynamics} and
one about \emph{architecture}, and our central claim is that they are
the same idea applied at two scales.

The first idea is that supervised learning on a linear model is already
a mechanical process. Given a data matrix $A \in \R^{m\times n}$ and a
target $b$, the gradient of the least--squares loss is
$\frac{1}{m} A^\top(Ax-b)$, and the transpose $A^\top$ is precisely the
operator that carries the residual from \emph{sample space} back into
\emph{feature space}, where the parameters live. If we additionally
endow feature space with a graph Laplacian $L_f$ built from $A^\top$ and
a mass matrix $M$, the resulting pair $(L_f, M)$ is a mass--spring
system in the sense of Rayleigh's \emph{Theory of
Sound}~\cite{rayleigh_theory_sound_1}: its generalised eigenvalues are
squared natural frequencies, its eigenvectors are mode shapes, and
gradient descent in this geometry behaves like a damped oscillator
settling in a potential well. This vibrational reading yields a
mass-aware Laplacian preconditioner with provable linear convergence and
a spectral (taumode) regulariser that biases learning toward the
energy-dispersion bands of a corpus.

The second idea concerns \emph{structured intermediate states} in
recurrent transformers. Transformers excel at flexible computation but
impose no constraint on intermediate representations: every hidden state
may occupy any point in $\R^d$, and structure must be learned entirely
from data. Recent work argues that forcing representations through a
constrained domain at each recurrent step is a strong inductive bias for
iterative reasoning. Lattice Deduction Transformers~\cite{davis2026ldt}
make this concrete for logical deduction: every recurrent step projects
its latent state onto a lattice of candidate sets, so inference becomes
a sequence of provably sound partial derivations.

We unify the two. The mass--spring system of the first idea \emph{is}
the constrained domain of the second. The eigenmodes of $L_f$ are the
``pure tones'' of the feature graph; their mixtures are spectral chords;
and a second-order damped wave step on this modal lattice is a structural
operator that constrains recurrent hidden states in exact analogy to
lattice projection in \textsc{LDT}. Where \textsc{LDT} projects onto a
logical lattice, the \VDT{} propagates along a \emph{spectral} lattice
governed by a physical wave equation.

\paragraph{Contributions.}
\begin{itemize}
  \item \textbf{Foundations (Sections~\ref{sec:dual_spaces}--\ref{sec:vibrational_gd}).}
        A self-contained mass-weighted vibrational interpretation of
        feature-space Laplacians: the dual-space role of $A^\top$, the
        pair $(L_f,M)$ and the generalised Rayleigh quotient $R_M$, a
        mass-aware Laplacian preconditioner $S_{\sigma,M}$ with a linear
        convergence guarantee, a taumode frequency regulariser, and the
        damped-oscillator reading of gradient descent that underlies the
        wave update.
  \item \textbf{Architecture (Section~\ref{sec:architecture}).}
        A \emph{vibrational deduction transformer} that replaces lattice
        projection with Laplacian modal propagation --- a second-order
        discrete wave update $\Phi_L$ with damping and forcing controlled
        by attention context --- as the structured intermediate operator
        in a recurrent transformer.
  \item \textbf{Signed density state (Section~\ref{sec:density_matrix}).}
        A signed density matrix $\varrho_t = \varrho_t^{+}-\varrho_t^{-}$
        in the Laplacian eigenbasis, supporting positive/negative
        evidential mass and interferometric cancellation, with
        measurable modal observables.
  \item \textbf{Objectives and protocol (Sections~\ref{sec:training}--\ref{sec:experiments}).}
        A depth-supervised training objective combining task loss,
        taumode frequency regularisation, and an occupancy penalty, and
        an evaluation protocol designed for exact comparability with the
        \textsc{LDT} benchmark, using modal energy spectra as
        interpretable diagnostics.
\end{itemize}

\paragraph{Organisation.}
Section~\ref{sec:dual_spaces} recalls the dual-space role of the
transpose. Section~\ref{sec:feature_graph} constructs the mass-weighted
feature graph and the generalised Rayleigh quotient.
Section~\ref{sec:cost} defines the Laplacian-regularised cost,
Section~\ref{sec:laplace_gd} the preconditioned gradient method and its
convergence, Section~\ref{sec:rayleigh_reg} the taumode regulariser, and
Section~\ref{sec:vibrational_gd} the vibrational (second-order) reading
of gradient descent. Sections~\ref{sec:architecture}--\ref{sec:experiments}
build the \VDT{} on these foundations. We close with discussion,
limitations, and conclusions.

%===================================================================
\part*{Part I --- Foundations: Vibrational Learning on Feature Space}
%===================================================================

%-------------------------------------------------------------------
\section{Dual spaces and the role of the transpose}
\label{sec:dual_spaces}
%-------------------------------------------------------------------

Let $V \cong \R^m$ be the sample space and $W \cong \R^n$ the feature
space. We interpret the data matrix $A \in \R^{m \times n}$ as a linear
map
\begin{equation}
    A : W \longrightarrow V, \quad x \mapsto A x.
\end{equation}
Under the standard Euclidean inner products on $V$ and $W$, the
transpose $A^\top$ is the adjoint map $A^\top : V \to W$, satisfying
\begin{equation}
    \langle A x, v \rangle_V = \langle x, A^\top v \rangle_W
    \quad \text{for all } x \in W, v \in V.
\end{equation}
Thus $A^\top$ transports vectors in sample space into feature space, and
can be viewed as a concrete realisation of the passage to the dual space
$V^*$.

Given a parameter vector $x \in W$, the residual is
$e(x) = A x - b \in V$, the least--squares cost is
$J(x) = \frac{1}{2m} \langle e(x), e(x) \rangle_V$, and its gradient is
\begin{equation}
    \nabla J(x)
    = \frac{1}{m} A^\top e(x)
    = \frac{1}{m} A^\top (A x - b).
\end{equation}
The factor $A^\top$ plays two roles. \emph{Dimensional alignment:} the
residual lives in $V \cong \R^m$ whereas the parameter vector lives in
$W \cong \R^n$, and $A^\top$ maps the residual into $W$, producing a
valid direction in parameter space. \emph{Error attribution:} each row
of $A^\top$ corresponds to a feature, and $A^\top e(x)$ computes the
correlation between each feature and the residual, attributing portions
of the error to coordinates of $x$. The normal equations
$A^\top (A x^\star - b) = 0$ then require that every feature functional
in this dual space sees no correlation with the residual. When $A$ is
rank-deficient or non-square the least-squares solution may be expressed
via the Moore--Penrose pseudoinverse
$A^+ = \lim_{\delta \to 0}(A^\top A + \delta I)^{-1}A^\top$;
see~\cite{golub_vanloan}.

%-------------------------------------------------------------------
\section{Feature graph and mass-weighted Rayleigh quotient}
\label{sec:feature_graph}
%-------------------------------------------------------------------

We endow the feature space $W$ with additional structure by constructing
a mass-weighted graph whose nodes are features and whose edges encode
feature--feature similarity. Crucially, this graph is built from the
transposed data $A^\top$.

Let $X = A^\top \in \R^{n \times m}$. The $i$-th row $X_{i:} \in \R^m$
is the signal of feature $i$ across all samples. We construct an
undirected weighted graph $G_f = (V_f, E_f, W_f)$ with node set
$V_f = \{1,\dots,n\}$ and symmetric weight matrix $W_f$ with entries
\begin{equation}
    (W_f)_{ij} = \kappa(X_{i:}, X_{j:}),
    \quad 1 \le i,j \le n,
\end{equation}
where $\kappa$ is a similarity kernel, for example a Gaussian kernel
$\kappa(u,v) = \exp(-\|u-v\|_2^2 / \sigma^2)$, a cosine similarity, or a
correlation measure. In practice $W_f$ is sparsified by keeping each
node's $k$ nearest neighbours. This data-driven construction of a
Laplacian from observed signals follows work on learning graph
Laplacians for smooth graph-signal representations~\cite{dong_etal_2016},
the key difference being that the ``signals'' are feature profiles across
samples rather than item-level observations.

\subsection{Laplacian, mass matrix, and Rayleigh quotient}

The degree matrix is $(D_f)_{ii} = \sum_{j} (W_f)_{ij}$ and the
combinatorial graph Laplacian is $L_f = D_f - W_f$. It is symmetric
positive semidefinite, and for any $z \in \R^n$,
\begin{equation}
    z^\top L_f z
    = \frac{1}{2} \sum_{i,j} (W_f)_{ij} (z_i - z_j)^2,
    \label{eq:lap_quadratic}
\end{equation}
which measures the smoothness of $z$ on the graph: it vanishes if and
only if $z$ is constant on each connected
component~\cite{merris_laplacian,spielman_laplacian}.

\begin{remark}[Combinatorial vs.\ normalised Laplacian]
We use the combinatorial Laplacian $L_f = D_f - W_f$ throughout. For
heterogeneous degrees the symmetric normalised Laplacian
$\mathcal{L} = D_f^{-1/2} L_f D_f^{-1/2}$ is often preferable: its
spectrum lies in $[0,2]$ regardless of graph scale, making $R_M$ more
comparable across graphs. With $M = D_f$, the generalised eigenproblem
$L_f u = \lambda D_f u$ is equivalent to the $\mathcal{L}$ eigenproblem,
so the two formulations coincide under the degree-mass choice.
\end{remark}

To obtain a vibrational system in the sense of Rayleigh we introduce a
positive definite diagonal mass matrix $M \in \R^{n \times n}$, with
$M_{ii} > 0$ assigning a mass to each feature.

\paragraph{Canonical constructions of $M$.}
\begin{itemize}
\item \emph{Degree mass} (default): $M = D_f$, free as a by-product of
      building $L_f$. This turns $L_f u = \lambda M u$ into the
      symmetric normalised Laplacian eigenproblem; highly connected
      features receive higher mass, matching the intuition that heavier
      oscillators are harder to move.
\item \emph{ArrowSpace taumode mass}: when an ArrowSpace index is
      available, set $M_{ii} = (1 - \lambda_i^\tau + \varepsilon)^{-1}$,
      where $\lambda_i^\tau \in [0,1)$ is the taumode score of feature
      $i$. Spectrally rough, high Rayleigh-energy features receive higher
      mass, coupling inertia to the ArrowSpace spectral geometry.
\end{itemize}

The pair $(L_f, M)$ defines a generalised eigenproblem
\begin{equation}
    L_f u_k = \lambda_k M u_k,
    \label{eq:eigen}
\end{equation}
whose eigenvalues $\lambda_k \ge 0$ are interpreted as squared natural
frequencies of the feature-space mass--spring system, and whose
eigenvectors $u_k$ are mode shapes. For any nonzero $z \in \R^n$ the
generalised Rayleigh quotient is
\begin{equation}
    R_M(z) = \frac{z^\top L_f z}{z^\top M z},
    \label{eq:rayleigh}
\end{equation}
playing exactly the role of Rayleigh's quotient in vibration analysis:
if $z$ aligns with an eigenmode $u_k$ then $R_M(z) = \lambda_k$, and more
generally $R_M(z)$ is an energy-weighted average of
eigenvalues~\cite{osher_lsgd_2022}. A natural Markov transition matrix on
the feature graph is
\begin{equation}
    P_f = D_f^{-1} W_f,
    \label{eq:markov}
\end{equation}
whose $i$-th row holds the transition probabilities of a random walk on
features. Thus $L_f$ and $M$ encode the geometry and inertia of feature
space induced by $A^\top$, while $P_f$ encodes the pathways along which
information flows between features. These three objects ---
$L_f$, $M$, and $P_f$ --- recur throughout the architecture of
Part~II.

\subsection{Relationship to ArrowSpace and Graph Wiring}
\label{sec:arrowspace}

The feature Laplacian $L_f$ and the taumode coordinates above are
exactly the constructions used by the ArrowSpace/Graph-Wiring spectral
search framework~\cite{moriondo_arrowspace_2025,moriondo_graphwiring_2026}.
There, a feature Laplacian built from a corpus matrix $A^\top$ is paired
with taumode values $\lambda_i^\tau \in [0,1)$ that compress the Rayleigh
energy of each corpus vector into a one-dimensional spectral coordinate
for vector search and retrieval. In this paper the same $L_f$, its
eigendecomposition $L_f = U\Lambda U^\top$, the mass matrix $M$, and the
prototype taumode values all serve a learning role: they define the
spectral lattice on which the vibrational dynamics of
Part~II propagate.

%-------------------------------------------------------------------
\section{A mass-weighted Laplacian-regularised cost}
\label{sec:cost}
%-------------------------------------------------------------------

Equipped with $L_f$ and $M$, we define a Laplacian-regularised
least--squares functional
\begin{equation}
    J_\lambda(x)
    = \frac{1}{2m} \|A x - b\|_2^2
      + \frac{\lambda}{2} x^\top L_f x,
    \quad x \in \R^n,
    \label{eq:laplacian_cost}
\end{equation}
with $\lambda \ge 0$. The first term enforces data fidelity; the second
penalises lack of smoothness of $x$ across the feature graph, so that
strongly connected features are encouraged to have similar coefficients,
mirroring classical graph-based semi-supervised
learning~\cite{zhu_ghahramani_lafferty_2003,belkin_matveeva_niyogi_2004,smola_kondor_2003}.
Using~\eqref{eq:rayleigh}, the smoothness term factors as
\begin{equation}
    x^\top L_f x = R_M(x)\, x^\top M x,
    \label{eq:rayleigh_factorisation}
\end{equation}
an interpretive lens reading the penalty as a mass-weighted squared
frequency $R_M(x)$ times a mass-weighted energy $x^\top M x$.

The gradient is
$\nabla J_\lambda(x) = \frac{1}{m} A^\top (A x - b) + \lambda L_f x$,
and the stationary condition is
\begin{equation}
    \left( \frac{1}{m} A^\top A + \lambda L_f \right) x
    = \frac{1}{m} A^\top b,
    \label{eq:laplacian_normal}
\end{equation}
so the Laplacian regulariser modifies the classical normal equations by
adding a graph-based smoothness term in feature space.

%-------------------------------------------------------------------
\section{Laplacian-preconditioned gradient descent}
\label{sec:laplace_gd}
%-------------------------------------------------------------------

We now describe a learning algorithm that uses the mass-weighted
feature-space Laplacian to define the pathways along which gradients are
diffused at each step.

\subsection{Explicit and implicit smoothing}

For the unregularised cost $J(x) = \frac{1}{2m}\|Ax-b\|_2^2$, gradient
descent updates $x_{t+1} = x_t - \eta \frac{1}{m} A^\top (Ax_t - b)$,
using $A^\top$ purely to move the residual into feature space. To
incorporate feature-graph pathways we smooth the gradient
$g_t = \nabla J(x_t)$. The \emph{explicit} smoothing is
\begin{equation}
    \tilde g_t = (I - \sigma L_f) g_t,
    \label{eq:explicit_diffusion}
\end{equation}
which is positive semidefinite only when $\sigma < 1/\lambda_{\max}(L_f)$;
violating this bound makes the operator indefinite and can diverge. The
\emph{implicit} mass-aware smoothing uses the resolvent
\begin{equation}
    S_{\sigma,M} = (M + \sigma L_f)^{-1} M,
    \qquad
    \tilde g_t = S_{\sigma,M} g_t,
    \label{eq:implicit_diffusion}
\end{equation}
which is positive semidefinite for \emph{all} $\sigma \ge 0$ and
$M \succ 0$, removing the stability restriction on $\sigma$ --- the main
reason it is preferred. $S_{\sigma,M}$ is one step of implicit Euler
discretisation of the mass-weighted heat equation
$M \dot z + L_f z = 0$ on the feature graph; setting $M = I$ recovers the
isotropic smoothing $(I + \sigma L_f)^{-1}$~\cite{osher_lsgd_2022}. The
implicit operator $S_{\sigma,M}$ reappears in Part~II as the stable
alternative to the explicit wave step.

\subsection{Convergence in the quadratic case}
\label{subsec:convergence}

Writing $H = \frac{1}{m} A^\top A$ and $P_{\sigma,M} = S_{\sigma,M}$, the
preconditioned update is
$x_{t+1} = x_t - \eta P_{\sigma,M} H (x_t - x^\star)$, where $x^\star$ is
the unique minimiser of $J$ whenever $H \succ 0$.

\begin{proposition}[Linear convergence under mass-aware smoothing]
\label{prop:convergence}
Assume $H \succ 0$ and $L_f \succeq 0$ symmetric, with $M \succ 0$
diagonal. For any $\sigma > 0$ the preconditioner
$P_{\sigma,M} = (M + \sigma L_f)^{-1} M$ is symmetric positive definite,
and the spectrum of $\mathcal{H}_{\sigma,M} = P_{\sigma,M} H$ lies in
$(0,\infty)$. Let
$\mu_{\sigma,M} = \lambda_{\min}(\mathcal{H}_{\sigma,M})$ and
$L_{\sigma,M} = \lambda_{\max}(\mathcal{H}_{\sigma,M})$. If
$0 < \eta < 2/L_{\sigma,M}$, then in the $P_{\sigma,M}^{-1}$-weighted
norm $\|w\|_{P^{-1}}^2 = w^\top P_{\sigma,M}^{-1} w$,
\[
    \|x_{t+1} - x^\star\|_{P^{-1}}
    \le (1 - \eta \mu_{\sigma,M})\,\|x_t - x^\star\|_{P^{-1}},
\]
and since $P_{\sigma,M} \succ 0$ makes this norm equivalent to the
Euclidean norm, the iterates converge linearly to $x^\star$.
\end{proposition}

\begin{proof}[Proof sketch]
\emph{(i)} Since $M \succ 0$ is diagonal and $L_f \succeq 0$,
$A_\sigma = M + \sigma L_f \succ 0$. Writing
$A_\sigma = M^{1/2}(I + \sigma M^{-1/2} L_f M^{-1/2})M^{1/2}$ shows
$P_{\sigma,M} = A_\sigma^{-1} M$ is congruent via $M^{1/2}$ to the SPD
matrix $(I + \sigma M^{-1/2} L_f M^{-1/2})^{-1}$, so
$P_{\sigma,M} \succ 0$.
\emph{(ii)} Because $P_{\sigma,M} \succ 0$ and $H \succ 0$,
$\mathcal{H}_{\sigma,M}$ shares its spectrum with the SPD matrix
$P_{\sigma,M}^{1/2} H P_{\sigma,M}^{1/2}$, hence all positive.
\emph{(iii)} The error obeys
$e_{t+1} = (I - \eta\mathcal{H}_{\sigma,M})e_t$; under
$\tilde e = P_{\sigma,M}^{-1/2} e$ the iteration matrix becomes the
symmetric $I - \eta P_{\sigma,M}^{1/2} H P_{\sigma,M}^{1/2}$, whose
eigenvalues lie in $[\mu_{\sigma,M},L_{\sigma,M}]\subset(0,\infty)$. For
$0<\eta<2/L_{\sigma,M}$ they all lie in $(-1,1)$, giving the stated
contraction; norm equivalence transfers it to $\|\cdot\|_2$. The full
argument is in Appendix~\ref{app:proof_sketch}.
\end{proof}

The spectrum of $\mathcal{H}_{\sigma,M}$ is fully characterised by the
eigenvalues of $S_{\sigma,M} H = (M + \sigma L_f)^{-1} M (A^\top A/m)$,
so the allowable step size and contraction factor are computable from
the preconditioned Hessian: Laplacian smoothing changes the geometry of
the level sets through $S_{\sigma,M}$, but the standard spectral theory
of gradient descent on quadratics continues to apply.
Proposition~\ref{prop:convergence} concerns the base convex model and is
independent of the non-convex taumode regulariser of
Section~\ref{sec:rayleigh_reg}.

\subsection{Algorithmic summary}

We summarise mass-aware Laplacian-preconditioned gradient descent as
Algorithm~\ref{alg:laplace_gd}.

\begin{algorithm}[ht]
\begin{center}
\fbox{\begin{minipage}{0.95\textwidth}
\textbf{Mass-Weighted Laplacian-Preconditioned Gradient Descent}\\[0.5ex]
\emph{Input:} $A \in \R^{m \times n}$, target $b \in \R^m$, step size
$\eta > 0$, diffusion $\sigma \ge 0$, iterations $T$.\\[0.5ex]
\emph{Preprocessing:}
\begin{enumerate}
\item Form $X = A^\top$. Build affinity $W_f$ from rows of $X$ via a
      kernel $\kappa$ and sparsification.
\item Compute $D_f$, $L_f = D_f - W_f$, and a positive diagonal mass
      matrix $M$ (degree mass or ArrowSpace taumode mass).
\item Optionally form $S_{\sigma,M} = (M + \sigma L_f)^{-1} M$.
\end{enumerate}
\emph{For} $t = 0,\dots,T-1$:
\begin{enumerate}
\item Residual $r_t = A x_t - b$; raw gradient
      $g_t = \frac{1}{m} A^\top r_t$.
\item Smooth: explicit $\tilde g_t = (I - \sigma L_f) g_t$ or implicit
      $\tilde g_t = S_{\sigma,M} g_t$.
\item Update $x_{t+1} = x_t - \eta \tilde g_t$.
\end{enumerate}
\emph{Output:} $x_T$.
\end{minipage}}
\end{center}
\caption{Mass-weighted Laplacian-preconditioned gradient descent.}
\label{alg:laplace_gd}
\end{algorithm}

%-------------------------------------------------------------------
\section{Taumode frequency regularisation}
\label{sec:rayleigh_reg}
%-------------------------------------------------------------------

\emph{This section is a speculative extension: the taumode regulariser
$J_{\mathrm{freq}}$ is a research direction rather than a validated
component, and Proposition~\ref{prop:convergence} is independent of it.}
The mass--spring system $(L_f,M)$ lets us go beyond smoothness penalties
and bias the frequencies of update directions toward the spectral bands
occupied by in-distribution corpus items in ArrowSpace.

At iteration $t$ consider trial shapes $q_j(x_t) \in \R^n$, e.g.\
$q_1 = \nabla J(x_t)$ and optionally Krylov directions
$q_2 = L_f \nabla J(x_t)$. Rather than computing an exact Rayleigh
quotient per trial shape at every step, we use precomputed taumode
values as a low-dimensional spectral proxy. Let
$\{\phi_c\}_{c=1}^C$ be feature-space prototypes (Graph Wiring
centroids) with taumode values $\lambda^\tau_c \in [0,1)$ computed
offline by ArrowSpace~\cite{moriondo_arrowspace_2025,moriondo_graphwiring_2026}.
Projecting $q_j(x_t)$ onto the prototypes with softmax similarity weights
$w_{j,c}(x_t)$ yields a differentiable approximate frequency
\begin{equation}
    \hat\rho_j(x_t) = \sum_{c=1}^C w_{j,c}(x_t)\,\lambda^\tau_c,
    \label{eq:approx_freq}
\end{equation}
a one-dimensional summary of the spectral position of $q_j(x_t)$
relative to the corpus taumode spectrum. From the empirical taumode
distribution we estimate a fundamental scale $\lambda_\star^\tau$ as a
robust central statistic (the median of the $25$--$75$th percentile
band), clamping $\hat\rho_j$ to this band so the regulariser acts chiefly
when frequencies leave the in-distribution range. The taumode frequency
MSE term and full functional are
\begin{equation}
    J_{\mathrm{freq}}(x)
    = \frac{1}{2J} \sum_{j=1}^J \bigl( \hat\rho_j(x) - \lambda_\star^\tau \bigr)^2,
    \qquad
    J_{\lambda,\mu}(x) = \frac{1}{2m}\|Ax-b\|_2^2
      + \frac{\lambda}{2} x^\top L_f x + \mu\, J_{\mathrm{freq}}(x).
    \label{eq:freq_cost}
\end{equation}
Since $\hat\rho_j$ depends nonlinearly on $x$, $J_{\mathrm{freq}}$ is
generally non-convex and is best viewed as a heuristic spectral bias;
the base quadratic-with-preconditioning model remains convex.
This taumode MSE reappears verbatim, lifted to the recurrent trajectory,
as the frequency loss $\mathcal{L}_{\rm freq}$ of the \VDT{} in
Section~\ref{sec:training}.

%-------------------------------------------------------------------
\section{Gradient descent as a settling vibrational system}
\label{sec:vibrational_gd}
%-------------------------------------------------------------------

The preconditioned method above is first-order. The deeper mechanical
content --- and the bridge to the architecture of Part~II --- comes from
treating learning as a \emph{second-order} damped vibration.

\subsection{Loss landscapes as vibrational potentials}

Treat the parameters $x \in \R^n$ as generalised coordinates of a
fictitious mechanical system and the loss $J(x)$ as its potential
energy. With kinetic energy $\frac{1}{2}\dot x^\top M \dot x$ and a
damping matrix $C \succeq 0$ we consider the damped second-order dynamics
\begin{equation}
    M \ddot x(t) + C \dot x(t) + \nabla J(x(t)) = 0,
    \label{eq:vibrational_gd_continuous}
\end{equation}
the analogue of Rayleigh's damped vibration with $\nabla J$ as the
restoring force. Near a non-degenerate minimum $x^\star$ with Hessian
$H = \nabla^2 J(x^\star) \succ 0$, writing $u = x - x^\star$,
\begin{equation}
    M \ddot u + C \dot u + H u = 0,
    \label{eq:linearized_vibration}
\end{equation}
exactly the form analysed by Rayleigh: $M$ is mass, $H$ is stiffness, and
the vibrational modes are governed by the eigenpairs of $M^{-1}H$, with
squared frequencies given by the generalised Rayleigh quotient
$R_M(z) = z^\top H z / z^\top M z$. Low $R_M(z)$ are slow modes of the
optimiser; high values are stiff directions.

\subsection{From overdamped descent to explicit vibration}

Standard gradient descent $x_{k+1} = x_k - \eta \nabla J(x_k)$ is the
explicit Euler discretisation of the first-order flow
$\dot x = -\nabla J(x)$, recovered from
\eqref{eq:vibrational_gd_continuous} in the overdamped limit
$M \to 0$, $C = I$. Thus ordinary descent explores only the zero-inertia
regime. Retaining inertia and discretising the full second-order system
with step $h$ (and $M=I$ for simplicity),
\begin{align}
    v_{k+1} &= v_k - h \nabla J(x_k) - \gamma h v_k, \label{eq:velocity_update}\\
    x_{k+1} &= x_k + h v_{k+1}, \label{eq:position_update}
\end{align}
a discrete analogue of $\ddot x + \gamma \dot x + \nabla J(x) = 0$. Along
each eigenvector $z$ of $H$ with $Hz = \lambda z$, the mode obeys
$\ddot u^\lambda + \gamma \dot u^\lambda + \lambda u^\lambda = 0$ --- a
damped harmonic oscillator with natural frequency $\sqrt{\lambda}$.
Rayleigh's classification applies directly: underdamped
($\gamma^2 < 4\lambda$, oscillatory decay), critical
($\gamma^2 = 4\lambda$, fastest non-oscillatory return), and overdamped
($\gamma^2 > 4\lambda$, slow monotone relaxation). Choosing $M$ to encode
feature masses and $C$ for mode-dependent damping yields a
spectrum-aware optimiser interpretable through Rayleigh's analysis.

\begin{remark}[The seed of the wave update]
\label{rem:seed}
Equations~\eqref{eq:velocity_update}--\eqref{eq:position_update} are the
scalar, single-parameter prototype of the matrix-valued vibrational
state update $\Phi_L$ that drives the \VDT{}
(Section~\ref{sec:vibrational_block}). There the stiffness $H$ is
replaced by the feature Laplacian $L_f$, the single coordinate $x$ by a
full token--feature state $Q_t$, and the constant forcing by an
attention-derived term $B_t$. The damped-oscillator modal decoupling of
this section is precisely what makes the architecture's spectral
formulation possible.
\end{remark}

%===================================================================
\part*{Part II --- Architecture: The Vibrational Deduction Transformer}
%===================================================================

%-------------------------------------------------------------------
\section{The Vibrational Deduction Transformer}
\label{sec:architecture}
%-------------------------------------------------------------------

We now lift the mechanics of Part~I from a single parameter vector to the
hidden states of a recurrent transformer. The mass--spring system
$(L_f,M)$ becomes the constrained domain through which every recurrent
hidden state must pass, playing the role that the logical lattice plays
in the Lattice Deduction Transformer~\cite{davis2026ldt}.

\subsection{Lattice Deduction Transformers as structural inspiration}

Davis et al.~\cite{davis2026ldt} introduce the \textsc{LDT}, a recurrent
transformer that approximates logically sound deduction by repeatedly
projecting its hidden state onto a lattice of candidate hypotheses. Each
recurrent step comprises a standard transformer layer and a lattice
projection $\Pi_\mathcal{L}$ mapping the latent state to the nearest
admissible lattice element. The structured intermediate state acts as a
powerful inductive bias: by constraining what the network can ``think''
at each step, the architecture converges faster and generalises better
on logical and arithmetic reasoning. We replace $\Pi_\mathcal{L}$ with
the Laplacian modal operator $\Phi_L$, and replace the logic lattice with
the spectral lattice of $L_f$ from Part~I.

\subsection{Overall recurrent structure}

Let $X_0 \in \R^{n \times d}$ be the input token features after
embedding, with $n$ tokens and feature dimension $d$. The \VDT{} iterates
a depth-$K$ recurrence
\begin{equation}
  X_{t+1} = \mathrm{TransformerBlock}\!\bigl(
    X_t + \Phi_L(X_t, Q_{t-1}, Q_t)
  \bigr),
  \quad t = 1,\dots,K,
  \label{eq:vdt_recurrence}
\end{equation}
where $\Phi_L$ is the vibrational state block
(Section~\ref{sec:vibrational_block}) and $Q_t \in \R^{n \times d}$ is
the vibrational state, an explicit second-order memory. We initialise
$Q_0 = Q_1 = X_0$ and update $Q_{t+1}$ via the discrete wave equation at
each step. The recurrent depth $K$ is a hyperparameter; all transformer
blocks share weights across recurrent steps, following the Universal
Transformer convention~\cite{dehghani2019universal} adopted also by
\textsc{LDT}.

\subsection{Feature graph construction}

Given $X_0 \in \R^{n \times d}$ we construct a feature graph $G_f$ over
the $d$ feature coordinates, exactly as in
Section~\ref{sec:feature_graph} but now per-batch and online. Let
$\tilde{X} = X_0^\top \in \R^{d \times n}$, whose $i$-th row is the
signal of feature $i$ across all $n$ tokens. The affinity is the
Gaussian kernel of Section~\ref{sec:feature_graph},
\begin{equation}
  (W_f)_{ij} = \exp\!\bigl(-\|\tilde{X}_{i:} - \tilde{X}_{j:}\|_2^2 / \sigma_f^2\bigr),
  \quad i \neq j;
  \quad (W_f)_{ii} = 0,
  \label{eq:affinity}
\end{equation}
sparsified by top-$k$ neighbours, with $D_f$, $L_f = D_f - W_f$, and
$P_f = D_f^{-1} W_f$ defined as before. For a sequence task, $G_f$ can be
fixed once per batch or updated per recurrent step at modest cost.

\subsection{Vibrational state block}
\label{sec:vibrational_block}

The vibrational state block implements a discrete second-order wave
equation on the feature graph --- the matrix-valued lift of the scalar
oscillator of Section~\ref{sec:vibrational_gd}
(Remark~\ref{rem:seed}). Let $Q_{t-1}, Q_t \in \R^{n \times d}$ be the
previous two states. The forcing term is derived from the current
transformer context $H_t = \mathrm{TransformerBlock}(X_t)$,
\begin{equation}
  B_t = \mathrm{Linear}(H_t) \in \R^{n \times d},
  \label{eq:forcing}
\end{equation}
and the vibrational update is
\begin{equation}
  Q_{t+1} = 2Q_t - Q_{t-1}
             - \Delta t^2 \, Q_t L_f^\top
             - \Gamma \odot (Q_t - Q_{t-1})
             + \Delta t^2 \, B_t,
  \label{eq:wave_update}
\end{equation}
where $\Delta t > 0$ is a learnable time-step (scalar or diagonal),
$\Gamma \in \R^{n \times d}$ is a learnable damping tensor initialised to
a small positive constant, $Q_t L_f^\top$ applies the Laplacian to each
token's feature vector (coupling feature coordinates through the graph),
and $\odot$ is the elementwise product. This is exactly the discretised
damped vibration $M\ddot x + C \dot x + \nabla J = 0$ of
\eqref{eq:vibrational_gd_continuous} with stiffness $L_f$ and forcing
$B_t$, applied to all $n$ tokens simultaneously: $L_f$ penalises
feature-direction differences across tokens according to edge weights,
and $\Gamma$ sets whether each mode is underdamped (oscillatory) or
overdamped (diffusive).

\begin{remark}[Stability]
The explicit step~\eqref{eq:wave_update} is stable when
$\Delta t^2 \lambda_{\max}(L_f) \leq 2$. In practice we clip the
learnable $\Delta t$ by $\Delta t \leq \sqrt{2/\lambda_{\max}(L_f)}$ at
each forward pass, or use a learnable implicit scheme analogous to the
mass-aware resolvent $S_{\sigma,M}$ of
Section~\ref{sec:laplace_gd}, which is unconditionally stable.
\end{remark}

\subsection{Spectral formulation}
\label{sec:spectral}

Let $L_f = U \Lambda U^\top$ be the eigendecomposition, $U$ orthogonal
and $\Lambda = \diag(\lambda_1,\dots,\lambda_d)$, $\lambda_k \ge 0$.
Define modal coordinates $\hat{Q}_t = Q_t U \in \R^{n \times d}$, so
$\hat{Q}_{t,k}$ projects all token features onto the $k$-th eigenmode. In
modal coordinates, \eqref{eq:wave_update} decouples mode-by-mode:
\begin{equation}
  \hat{Q}_{t+1,k} = 2\hat{Q}_{t,k} - \hat{Q}_{t-1,k}
    - \Delta t^2 \lambda_k \hat{Q}_{t,k}
    - \Gamma_k (\hat{Q}_{t,k} - \hat{Q}_{t-1,k})
    + \Delta t^2 \hat{B}_{t,k},
  \label{eq:modal_wave}
\end{equation}
a set of $d$ independent damped harmonic oscillators, one per mode, with
natural frequency $\omega_k = \sqrt{\lambda_k}$. This is exactly the
modal analysis of Section~\ref{sec:vibrational_gd}, now indexed by the
eigenmodes of $L_f$ rather than of a loss Hessian, and it is Rayleigh's
classical modal decomposition~\cite{rayleigh_theory_sound_1}.
Low-eigenvalue modes ($\lambda_k \approx 0$) are slowly varying global
features; high-eigenvalue modes are rapidly varying local features.
Learning adjusts $\Delta t$, $\Gamma_k$, and $\hat{B}_{t,k}$ to
distribute representational capacity across the modal spectrum.

%-------------------------------------------------------------------
\section{Signed Density Matrix State}
\label{sec:density_matrix}
%-------------------------------------------------------------------

\subsection{Motivation}

The vibrational state $Q_t$ encodes feature-space dynamics but treats all
modal occupancy as positive. In reasoning tasks it is natural to maintain
both \emph{supporting} and \emph{contradicting} evidence for a hypothesis
simultaneously, combining them through interference rather than simple
accumulation. We formalise this with a signed density matrix, which
generalises the positive-only Laplacian state of Part~I.

\subsection{Definition}

\begin{definition}[Signed density matrix state]
  The \emph{signed density matrix state} at step $t$ is a pair
  $(\varrho_t^{+}, \varrho_t^{-})$ with
  $\varrho_t^{+}, \varrho_t^{-} \in \R^{m \times m}$, $m \le d$ a chosen
  modal rank, both symmetric positive semidefinite, and the signed state
  \begin{equation}
    \varrho_t = \varrho_t^{+} - \varrho_t^{-}.
    \label{eq:signed_rho}
  \end{equation}
  Both components live in the Laplacian eigenbasis, i.e.\ the subspace
  spanned by the $m$ lowest-energy eigenvectors
  $U_m = [u_1|\cdots|u_m]$ of $L_f$.
\end{definition}

The diagonal entries $\varrho_{t,kk}^{\pm}$ are the positive/negative
occupancies of mode $k$, while off-diagonals capture multi-mode
coherences.

\subsection{Update equations}

Given the transformer context $H_t$, the projected context matrix
$C_t = U_m^\top H_t^\top H_t U_m \in \R^{m \times m}$ is automatically
PSD. Split it through learned affine maps,
\begin{equation}
  C_t^{+} = W_+^\top C_t W_+, \quad
  C_t^{-} = W_-^\top C_t W_-,
  \quad W_+, W_- \in \R^{m \times m},
  \label{eq:context_split}
\end{equation}
and update
\begin{align}
  \varrho_{t+1}^{+} &= (1-\alpha)\varrho_t^{+}
    + \alpha \, \mathrm{softplus}\bigl(C_t^{+}\bigr)
    + \beta \, \varrho_t^{+} \odot P_f,
    \label{eq:rho_plus} \\
  \varrho_{t+1}^{-} &= (1-\alpha)\varrho_t^{-}
    + \alpha \, \mathrm{softplus}\bigl(C_t^{-}\bigr)
    + \beta \, \varrho_t^{-} \odot P_f,
    \label{eq:rho_minus}
\end{align}
where $\alpha, \beta \in (0,1)$ are learnable, $P_f = D_f^{-1} W_f$ is the
graph Markov matrix of \eqref{eq:markov} (diffusion on feature modes),
and $\mathrm{softplus}(Z) = \log(1+\exp(Z))$ applied entrywise (then
symmetrised) preserves PSD structure. The diffusion term ties the
density-matrix dynamics to the same feature-graph pathways $P_f$ that
govern the preconditioner and wave update.

\subsection{Observables}

From the signed density matrix we extract scalar observables used as
diagnostics and as inputs to the prediction head:
\begin{align}
  E_t^{\pm} &= \Tr(\Lambda_m \varrho_t^{\pm})
    && \text{(modal energy of positive/negative components)}, \\
  C_t^{\rm coh} &= \Tr\bigl((\varrho_t^{+} - \varrho_t^{-})^2\bigr)
    && \text{(total coherence, squared Frobenius)}, \\
  \mathrm{occ}_{t,k} &= \varrho_{t,kk}^{+} - \varrho_{t,kk}^{-}
    && \text{(signed modal occupancy per mode $k$).}
\end{align}
Inspecting the modal energy spectrum during training reveals which
frequency bands the network uses most --- the learning-time analogue of
the taumode spectral coordinates of Section~\ref{sec:rayleigh_reg}.

\subsection{Reprojection into the \VDT{} recurrence}

The vibrational state is recovered from the density matrix by mixing the
transformer context with the signed modal structure,
\begin{equation}
  Q_t = H_t \, U_m \, \varrho_t \, U_m^\top + H_t,
  \label{eq:reprojection}
\end{equation}
before feeding back into \eqref{eq:vdt_recurrence}. This reprojection is
the \VDT{} analogue of \textsc{LDT}'s lattice projection
$\Pi_\mathcal{L}$: it constrains $Q_t$ to respect the spectral geometry
of $L_f$.

%-------------------------------------------------------------------
\section{Training Objectives}
\label{sec:training}
%-------------------------------------------------------------------

The full \VDT{} training objective is
\begin{equation}
  \mathcal{L} = \mathcal{L}_{\rm task}
              + \mu_1 \mathcal{L}_{\rm freq}
              + \mu_2 \mathcal{L}_{\rm depth},
  \label{eq:total_loss}
\end{equation}
where:
\begin{itemize}
  \item $\mathcal{L}_{\rm task}$ is the task-specific cross-entropy or
        MSE loss. Following \textsc{LDT}~\cite{davis2026ldt}, we supervise
        every intermediate step $t = 1,\dots,K$ with the ground-truth
        label and average. This \emph{depth supervision} forces the
        vibrational dynamics to represent useful partial conclusions at
        every recurrent depth.
  \item $\mathcal{L}_{\rm freq}$ is the taumode frequency MSE of
        Section~\ref{sec:rayleigh_reg}, lifted to the recurrent
        trajectory:
        \begin{equation}
          \mathcal{L}_{\rm freq}
          = \frac{1}{K} \sum_{t=1}^K
            \bigl(\hat\rho(Q_t) - \lambda_\star^\tau \bigr)^2,
          \label{eq:freq_loss}
        \end{equation}
        biasing the vibrational trajectory toward the in-distribution
        spectral band of the ArrowSpace corpus.
  \item $\mathcal{L}_{\rm depth}$ penalises excessive signed occupancy to
        prevent the density matrix from collapsing to a trivial state:
        \begin{equation}
          \mathcal{L}_{\rm depth}
          = \frac{1}{K} \sum_{t=1}^K \bigl\| \varrho_t \bigr\|_F^2 .
          \label{eq:depth_loss}
        \end{equation}
\end{itemize}
The objective thus inherits, term by term, the data-fidelity, spectral,
and stability ingredients of the Part~I cost functional
\eqref{eq:freq_cost}, now expressed over a recurrent forward-only
trajectory rather than a single least--squares solve.

%-------------------------------------------------------------------
\section{Experiments}
\label{sec:experiments}
%-------------------------------------------------------------------

\emph{This section presents the planned experimental protocol; results
will be reported in a subsequent revision. The protocol is designed for
exact comparability with Davis et al.~(2026)~\cite{davis2026ldt}.}

\subsection{Experimental protocol following LDT}

We adopt the three task families of Davis et al.:
\begin{enumerate}
  \item \textbf{Propositional satisfiability (SAT):} random 3-SAT
        instances of increasing clause-to-variable ratio; binary
        SAT/UNSAT label. Train up to 20 variables, test up to 40.
  \item \textbf{Syllogistic deduction:} multi-step categorical syllogisms
        (All $A$ are $B$; All $B$ are $C$; therefore All $A$ are $C$) with
        up to 8 steps, tested on chain lengths unseen in training.
  \item \textbf{Multi-step numeric reasoning:} integer arithmetic chains
        (e.g.\ iterative modular addition) of variable depth; accuracy
        vs.\ chain length.
\end{enumerate}

For each task we run four architectures under identical training budgets
(parameter count, compute steps, random seed sets):

\begin{center}
\begin{tabular}{ll}
\toprule
\textbf{Baseline} & \textbf{Description} \\
\midrule
Transformer & Vanilla recurrent transformer (shared weights, $K$ steps) \\
\textsc{LDT} & Lattice Deduction Transformer~\cite{davis2026ldt} \\
\VDT{} (ours) & This paper: wave update + density matrix \\
\VDT{}-ablation & \VDT{} without density matrix (plain $\Phi_L$ only) \\
\bottomrule
\end{tabular}
\end{center}

\subsection{Metrics}

For each task and architecture we report: accuracy vs.\ recurrent depth
$K \in \{2,4,8,16\}$; accuracy vs.\ problem complexity (clause ratio /
chain length / number of steps), including out-of-distribution lengths;
modal energy spectra $E_t^+, E_t^-$ as a function of depth (\VDT{} only);
and signed modal occupancy profiles
$\{\mathrm{occ}_{t,k}\}_{k=1}^m$ at convergence.

\subsection{Hypotheses}

From the theoretical properties of $\Phi_L$ and the signed density
matrix we hypothesise: (i) \VDT{} matches or exceeds \textsc{LDT} on
syllogistic deduction, via the correspondence between modal low-energy
states and globally consistent conclusions; (ii) \VDT{} generalises to
longer chains better than a plain recurrent transformer, because the
spectral constraint prevents representation collapse at large $K$;
(iii) the signed density matrix helps on SAT (where contradictions
matter) relative to the ablation; and (iv) modal energy spectra are
interpretable signatures --- SAT instances populate high-$\lambda$ modes
while definite syllogisms concentrate weight on low-$\lambda$ modes.

%-------------------------------------------------------------------
\section{Discussion}
\label{sec:discussion}
%-------------------------------------------------------------------

The \VDT{} uses the feature Laplacian as both \emph{geometry} (the
spectral lattice of Part~I) and \emph{dynamics} (the constrained wave
update of Part~II). Replacing \textsc{LDT}'s logical lattice with a
physical wave equation has three practical consequences.

\paragraph{Interpretability.}
Modal energies and occupancies are measurable during inference. For a
given input one can inspect which frequency bands are activated and
whether the signed density matrix has resolved (positive $\gg$ negative)
or remains ambiguous (balanced) --- a direct analogue of inspecting which
atoms in a lattice are activated in \textsc{LDT}, and a recurrent
counterpart to the taumode diagnostics of Part~I.

\paragraph{Generalisation by spectral regularity.}
The Laplacian constraint keeps high-depth representations spectrally
smooth; they cannot degenerate into arbitrary high-frequency noise. This
is the recurrent analogue of the convergence-stabilising role played by
the mass-aware preconditioner $S_{\sigma,M}$ in
Proposition~\ref{prop:convergence}, and it mirrors lattice projection's
role in keeping \textsc{LDT} states within logically consistent partial
derivations.

\paragraph{Connections to quantum machine learning.}
The signed density matrix $\varrho_t = \varrho_t^+ - \varrho_t^-$ is
formally similar to a difference of quantum states, though we make no
quantum claims. The analogy suggests techniques from open quantum systems
(Lindblad operators, decoherence) as alternative formulations of the
density-matrix update~\eqref{eq:rho_plus}--\eqref{eq:rho_minus}.

\paragraph{Limitations and future work.}
The explicit wave step~\eqref{eq:wave_update} imposes a stability
constraint on $\Delta t$ (mitigated by the implicit resolvent of
Part~I). The density-matrix update doubles per-step state variables. The
feature graph $G_f$ must be built or updated at training time; for large
$d$ this needs efficient sparse Laplacian solvers, for which the
ArrowSpace machinery is directly applicable. The taumode regulariser of
Section~\ref{sec:rayleigh_reg} remains speculative, and full
experimental results are pending. A formal stability characterisation of
taumode coordinates across corpora, and the interaction between
$\mathcal{L}_{\rm freq}$ and tau-modulated retrieval, are natural next
steps.

%-------------------------------------------------------------------
\section{Conclusion}
\label{sec:conclusion}
%-------------------------------------------------------------------

We have presented a single, self-contained account that runs from the
mechanics of least--squares to a recurrent reasoning architecture. The
foundations establish that feature space, equipped with a Laplacian
stiffness $L_f$ and mass $M$ built from $A^\top$, is a Rayleigh
mass--spring system whose generalised Rayleigh quotient governs
preconditioning (with provable linear convergence), spectral
regularisation, and the damped-oscillator behaviour of gradient descent.
Lifting this mechanism from a single parameter vector to the hidden
states of a recurrent transformer yields the Vibrational Deduction
Transformer, in which a second-order discrete wave equation and a signed
density matrix in the Laplacian eigenbasis constrain every forward step
to the spectral geometry of $L_f$ --- the structural analogue of lattice
projection in the Lattice Deduction Transformer. The result is an
architecture with an explicit mechanical pedigree, interpretable modal
diagnostics, and a benchmark protocol for direct comparison against
lattice-projection baselines; experimental validation is in progress.

\section*{Acknowledgements}
This work made use of LLM tools to assist with drafting and editing; all
technical claims remain the author's responsibility. Thanks to
\href{https://www.enverge.ai/}{Enverge Lab} for compute access and to
GitHub Sponsors for continued support; sponsor the project at
\href{https://github.com/sponsors/tuned-org-uk}{github.com/sponsors/tuned-org-uk}.

\begin{thebibliography}{99}

\bibitem{rayleigh_theory_sound_1}
Lord Rayleigh.
\newblock {\em The Theory of Sound, Volume I}, 2nd edition.
\newblock Macmillan, 1894.

\bibitem{wei2022cot}
J.~Wei, X.~Wang, D.~Schuurmans, M.~Bosma, B.~Ichter, F.~Xia, E.~Chi,
Q.~Le, and D.~Zhou.
\newblock Chain-of-thought prompting elicits reasoning in large language models.
\newblock In {\em Proc.\ NeurIPS}, 2022.

\bibitem{snell2024scaling}
C.~Snell, J.~Lee, K.~Xu, and A.~Kumar.
\newblock Scaling {LLM} test-time compute optimally can be more effective than
scaling model parameters.
\newblock arXiv:2408.03314, 2024.

\bibitem{austin2021d3pm}
J.~Austin, D.~D. Johnson, J.~Ho, D.~Tarlow, and R.~van den Berg.
\newblock Structured denoising diffusion models in discrete state spaces.
\newblock In {\em Proc.\ NeurIPS}, 2021.

\bibitem{golub_vanloan}
G.~H. Golub and C.~F. Van~Loan.
\newblock {\em Matrix Computations}, 4th edition.
\newblock Johns Hopkins University Press, 2013.

\bibitem{merris_laplacian}
R.~Merris.
\newblock Laplacian matrices of graphs: a survey.
\newblock {\em Linear Algebra and its Applications}, 197--198:143--176, 1994.

\bibitem{spielman_laplacian}
D.~A. Spielman.
\newblock Algorithms, graph theory, and linear equations in {L}aplacian matrices.
\newblock In {\em Proc.\ ICM}, 2010.

\bibitem{zhu_ghahramani_lafferty_2003}
X.~Zhu, Z.~Ghahramani, and J.~Lafferty.
\newblock Semi-supervised learning using Gaussian fields and harmonic functions.
\newblock In {\em Proc.\ ICML}, pages 912--919, 2003.

\bibitem{belkin_matveeva_niyogi_2004}
M.~Belkin, I.~Matveeva, and P.~Niyogi.
\newblock Regularization and semi-supervised learning on large graphs.
\newblock In {\em Proc.\ COLT}, pages 624--638, 2004.

\bibitem{smola_kondor_2003}
A.~J. Smola and R.~Kondor.
\newblock Kernels and regularization on graphs.
\newblock In {\em Learning Theory and Kernel Machines}, pages 144--158.
\newblock Springer, 2003.

\bibitem{dong_etal_2016}
X.~Dong, D.~Thanou, M.~Rabbat, and P.~Frossard.
\newblock Learning {L}aplacian matrix in smooth graph signal representations.
\newblock {\em IEEE Trans.\ Signal Processing}, 64(23):6160--6173, 2016.

\bibitem{osher_lsgd_2022}
S.~Osher, W.~Shi, and W.~Zhu.
\newblock {L}aplacian smoothing gradient descent.
\newblock {\em Research in the Mathematical Sciences}, 9(3):Paper No.~39, 2022.

\bibitem{moriondo_arrowspace_2025}
L.~Moriondo.
\newblock ArrowSpace: Spectral indexing for vector search.
\newblock {\em Journal of Open Source Software}, 10(90):9002, 2025.

\bibitem{moriondo_graphwiring_2026}
L.~Moriondo.
\newblock Graph Wiring and ArrowSpace: Spectral pathways for structure-aware learning.
\newblock Technical report, 2026.

\bibitem{davis2026ldt}
L.~Davis, M.~Jones, S.~Kumar, and P.~Williams.
\newblock Lattice Deduction Transformers.
\newblock arXiv:2605.08605v1, 2026.

\bibitem{dehghani2019universal}
M.~Dehghani, S.~Gouws, O.~Vinyals, J.~Uszkoreit, and {\L}.~Kaiser.
\newblock Universal Transformers.
\newblock In {\em Proc.\ ICLR}, 2019.

\bibitem{brown2024large}
OpenAI.
\newblock {OpenAI o1} System Card.
\newblock Technical report, OpenAI, September 2024.
\newblock \url{https://openai.com/index/openai-o1-system-card/}

\bibitem{muennighoff2025scaling}
OpenAI.
\newblock {OpenAI o3} and {o4-mini} System Card.
\newblock Technical report, OpenAI, April 2025.
\newblock \url{https://openai.com/index/o3-and-o4-mini-system-card/}

\end{thebibliography}

\newpage
\appendix

%-------------------------------------------------------------------
\section{Proof of Proposition~\ref{prop:convergence}}
\label{app:proof_sketch}
%-------------------------------------------------------------------

This appendix expands the proof of Proposition~\ref{prop:convergence},
the linear-convergence guarantee for mass-aware
Laplacian-preconditioned gradient descent on
$J(x) = \frac{1}{2m}\|A x - b\|_2^2$, with
$H = \tfrac{1}{m} A^\top A \succ 0$, $L_f \succeq 0$ symmetric,
$M \succ 0$ diagonal, and
$P_{\sigma,M} = (M + \sigma L_f)^{-1} M$, $\sigma > 0$.

\subsection{Step 1: $P_{\sigma,M}$ is symmetric positive definite}

Write $A_\sigma = M + \sigma L_f$. Since $M \succ 0$ diagonal and
$L_f \succeq 0$, $A_\sigma \succ 0$ for every $\sigma > 0$, so
$A_\sigma^{-1}$ exists. Factor through the symmetric square root of $M$,
\[
    A_\sigma
    = M^{1/2}\bigl(I + \sigma\, M^{-1/2} L_f M^{-1/2}\bigr) M^{1/2}
    =: M^{1/2}\,\widetilde A_\sigma\, M^{1/2},
\]
with $\widetilde A_\sigma = I + \sigma \widetilde L_f$ and
$\widetilde L_f = M^{-1/2} L_f M^{-1/2} \succeq 0$. Then
$P_{\sigma,M} = A_\sigma^{-1} M = M^{-1/2} \widetilde A_\sigma^{-1} M^{1/2}$
is congruent via $M^{1/2}$ to the SPD matrix
$\widetilde A_\sigma^{-1} = (I + \sigma \widetilde L_f)^{-1}$, confirming
$P_{\sigma,M} \succ 0$ and the existence of $P_{\sigma,M}^{1/2}$.

\subsection{Step 2: positive spectrum of the preconditioned Hessian}

The preconditioned Hessian $\mathcal{H}_{\sigma,M} = P_{\sigma,M} H$ is
similar to a symmetric matrix:
\[
    P_{\sigma,M}^{-1/2}\,\mathcal{H}_{\sigma,M}\, P_{\sigma,M}^{1/2}
    = P_{\sigma,M}^{1/2} H P_{\sigma,M}^{1/2}
    =: \mathcal{S}_{\sigma,M},
\]
which is SPD because $H \succ 0$ and $P_{\sigma,M}^{1/2} \succ 0$. Hence
$\mathcal{H}_{\sigma,M}$ and $\mathcal{S}_{\sigma,M}$ share the same real,
positive spectrum, and we set
$\mu_{\sigma,M} = \lambda_{\min}(\mathcal{S}_{\sigma,M}) > 0$,
$L_{\sigma,M} = \lambda_{\max}(\mathcal{S}_{\sigma,M}) < \infty$.

\subsection{Step 3: contraction in the weighted norm}

The error $e_t = x_t - x^\star$ obeys
$e_{t+1} = (I - \eta\,\mathcal{H}_{\sigma,M})\, e_t$. Because
$\mathcal{H}_{\sigma,M}$ is not symmetric in general, we measure the
error in the $P_{\sigma,M}^{-1}$-weighted norm
$\|w\|_{P^{-1}}^2 = w^\top P_{\sigma,M}^{-1} w$. Under
$\tilde e_t = P_{\sigma,M}^{-1/2} e_t$,
\[
    \tilde e_{t+1} = \bigl(I - \eta\, \mathcal{S}_{\sigma,M}\bigr)\,\tilde e_t,
\]
with $I - \eta\,\mathcal{S}_{\sigma,M}$ symmetric and eigenvalues
$1 - \eta\lambda$ for $\lambda \in [\mu_{\sigma,M}, L_{\sigma,M}]$. For
$0 < \eta < 2/L_{\sigma,M}$ every such eigenvalue lies in $(-1,1)$, so
\[
    \|I - \eta\,\mathcal{S}_{\sigma,M}\|_2
    = \max\bigl\{|1-\eta\mu_{\sigma,M}|,\,|1-\eta L_{\sigma,M}|\bigr\}
    \le 1 - \eta\mu_{\sigma,M} < 1.
\]
Since $\|e_t\|_{P^{-1}} = \|\tilde e_t\|_2$,
$\|x_{t+1} - x^\star\|_{P^{-1}} \le (1 - \eta\mu_{\sigma,M})\,\|x_t - x^\star\|_{P^{-1}}$.

\subsection{Step 4: equivalence of norms}

Because $P_{\sigma,M} \succ 0$,
\[
    \lambda_{\max}(P_{\sigma,M})^{-1/2}\,\|w\|_2
    \le \|w\|_{P^{-1}}
    \le \lambda_{\min}(P_{\sigma,M})^{-1/2}\,\|w\|_2 ,
\]
so iterating the contraction gives linear convergence in $\|\cdot\|_2$,
\[
    \|x_t - x^\star\|_2
    \le \sqrt{\kappa(P_{\sigma,M})}\,
         (1 - \eta\mu_{\sigma,M})^t\,
         \|x_0 - x^\star\|_2 ,
\]
with $\kappa(P_{\sigma,M}) = \lambda_{\max}(P_{\sigma,M})/\lambda_{\min}(P_{\sigma,M})$.
This completes the proof.
\hfill$\qed$

%-------------------------------------------------------------------
\section{PyTorch Implementation of the VDT Core}
\label{app:pytorch_vdt}
%-------------------------------------------------------------------

\subsection{Feature graph and Laplacian}

\begin{lstlisting}[language=Python,
  caption={Build the feature-space Laplacian from a token embedding matrix.}]
import torch
import torch.nn.functional as F

def build_feature_laplacian(X, k=5, sigma=None):
    """
    X : (n_tokens, d_features)
    Returns L_f (d x d), W_f (d x d), eigenvalues, eigenvectors.
    """
    Xt = X.T                             # (d, n)
    diff = Xt.unsqueeze(1) - Xt.unsqueeze(0)   # (d, d, n)
    dist2 = (diff ** 2).sum(-1)          # (d, d)

    if sigma is None:
        sigma = dist2.mean().sqrt().item() + 1e-8

    W = torch.exp(-dist2 / sigma**2)
    W.fill_diagonal_(0.0)

    # k-NN sparsification
    topk_vals, _ = torch.topk(W, k=k, dim=-1)
    threshold = topk_vals[:, -1:]
    W = W * (W >= threshold).float()
    W = (W + W.T) / 2           # symmetrise

    d_vec = W.sum(dim=1)        # degree vector
    D = torch.diag(d_vec)
    L = D - W

    # Eigendecomposition
    eigvals, eigvecs = torch.linalg.eigh(L)   # sorted ascending
    return L, W, eigvals, eigvecs
\end{lstlisting}

\subsection{Vibrational state block}

\begin{lstlisting}[language=Python,
  caption={Vibrational state block implementing the discrete wave update.}]
import torch
import torch.nn as nn

class VibrationalStateBlock(nn.Module):
    def __init__(self, d_feat, lambda_max, dt_init=0.1):
        super().__init__()
        self.d = d_feat
        self.lambda_max = lambda_max
        # Learnable time-step, clipped for stability
        self.log_dt = nn.Parameter(torch.tensor(dt_init).log())
        # Learnable per-feature damping
        self.log_gamma = nn.Parameter(
            torch.full((d_feat,), -2.0)    # init gamma ~ 0.135
        )
        # Forcing projection
        self.forcing = nn.Linear(d_feat, d_feat, bias=False)

    def forward(self, Q_prev, Q_curr, H_t, L_f):
        """
        Q_prev, Q_curr : (n, d)  vibrational states at t-1 and t
        H_t            : (n, d)  transformer context at t
        L_f            : (d, d)  feature Laplacian
        Returns Q_next : (n, d)
        """
        # Clip dt for CFL stability
        dt_max = (2.0 / (self.lambda_max + 1e-8)) ** 0.5
        dt = torch.exp(self.log_dt).clamp(max=dt_max)
        gamma = torch.exp(self.log_gamma)          # (d,)
        B_t = self.forcing(H_t)                    # (n, d)

        # Wave update: Q_{t+1} = 2Q - Q_prev - dt^2 Q L^T
        #                        - gamma*(Q - Q_prev) + dt^2 B
        stiffness = dt**2 * (Q_curr @ L_f.T)       # (n, d)
        damp = gamma.unsqueeze(0) * (Q_curr - Q_prev)  # (n, d)
        Q_next = (2 * Q_curr - Q_prev
                  - stiffness
                  - damp
                  + dt**2 * B_t)
        return Q_next
\end{lstlisting}

\subsection{Signed density matrix}

\begin{lstlisting}[language=Python,
  caption={Signed density-matrix state in the Laplacian eigenbasis.}]
class SignedDensityMatrix(nn.Module):
    def __init__(self, m_modes):
        """m_modes: number of low-energy modes to retain."""
        super().__init__()
        self.m = m_modes
        self.W_plus  = nn.Parameter(torch.eye(m_modes))
        self.W_minus = nn.Parameter(torch.eye(m_modes) * 0.1)
        self.alpha   = nn.Parameter(torch.tensor(0.3))
        self.beta    = nn.Parameter(torch.tensor(0.1))

    def forward(self, rho_plus, rho_minus, H_t, U_m, P_f_m):
        """
        rho_plus, rho_minus : (m, m)  current PSD states
        H_t     : (n, d)  transformer context
        U_m     : (d, m)  first m eigenvectors of L_f
        P_f_m   : (m, m)  P_f restricted to m modes
        Returns updated rho_plus, rho_minus : (m, m)
        """
        alpha = self.alpha.sigmoid()
        beta  = self.beta.sigmoid()

        # Context projected to modal subspace
        C = U_m.T @ H_t.T @ H_t @ U_m   # (m, m)

        # Positive / negative context projections
        C_plus  = self.W_plus.T  @ C @ self.W_plus
        C_minus = self.W_minus.T @ C @ self.W_minus

        # Ensure PSD via softplus (element-wise, then symmetrise)
        def make_psd(Z):
            Z = F.softplus(Z)
            return (Z + Z.T) / 2

        rho_plus_new  = ((1-alpha)*rho_plus
                         + alpha*make_psd(C_plus)
                         + beta*(rho_plus * P_f_m))
        rho_minus_new = ((1-alpha)*rho_minus
                         + alpha*make_psd(C_minus)
                         + beta*(rho_minus * P_f_m))
        return rho_plus_new, rho_minus_new

    def to_vibrational_state(self, rho_plus, rho_minus,
                              H_t, U_m):
        """Reprojection: mix signed density matrix into transformer context."""
        rho = rho_plus - rho_minus                  # (m, m)
        modal_correction = H_t @ U_m @ rho @ U_m.T  # (n, d)
        return H_t + modal_correction
\end{lstlisting}

\subsection{Full VDT forward pass}

\begin{lstlisting}[language=Python,
  caption={Sketch of the full VDT recurrent forward pass.}]
class VDT(nn.Module):
    def __init__(self, d, m_modes, K, n_heads=4, lambda_max=1.0):
        super().__init__()
        self.K = K
        self.transformer = nn.TransformerEncoderLayer(
            d_model=d, nhead=n_heads, batch_first=True
        )
        self.vib_block = VibrationalStateBlock(d, lambda_max)
        self.density   = SignedDensityMatrix(m_modes)
        self.m = m_modes

    def forward(self, X0, L_f, eigvecs, P_f):
        """
        X0     : (batch, n, d)
        L_f    : (d, d)
        eigvecs: (d, d)  columns are eigenvectors of L_f
        P_f    : (d, d)  Markov matrix on feature graph
        """
        B, n, d = X0.shape
        U_m = eigvecs[:, :self.m]    # (d, m) low-energy modes
        P_m = U_m.T @ P_f @ U_m     # (m, m)

        Q_prev = X0.clone()
        Q_curr = X0.clone()
        rho_p  = torch.zeros(self.m, self.m)
        rho_m  = torch.zeros(self.m, self.m)

        logits_all = []
        for t in range(self.K):
            # 1. Transformer mixing
            H_t = self.transformer(Q_curr)       # (B, n, d)

            # 2. Vibrational update (batch dimension: loop or vmap)
            Q_next_list = []
            for b in range(B):
                Qn = self.vib_block(Q_prev[b], Q_curr[b],
                                     H_t[b], L_f)
                Q_next_list.append(Qn)
            Q_next = torch.stack(Q_next_list)    # (B, n, d)

            # 3. Density matrix update (averaged over batch)
            H_avg = H_t.mean(0)                  # (n, d)
            rho_p, rho_m = self.density(
                rho_p, rho_m, H_avg, U_m, P_m
            )

            # 4. Reprojection
            X_repr_list = []
            for b in range(B):
                Xr = self.density.to_vibrational_state(
                    rho_p, rho_m, Q_next[b], U_m
                )
                X_repr_list.append(Xr)
            X_repr = torch.stack(X_repr_list)    # (B, n, d)

            # Collect per-depth logits for depth supervision
            logits_all.append(X_repr[:, 0, :])   # CLS token

            Q_prev = Q_curr
            Q_curr = X_repr

        return Q_curr, logits_all, (rho_p, rho_m)
\end{lstlisting}

%-------------------------------------------------------------------
\section{LDT-Mirrored Benchmark Tasks}
\label{app:ldt_tasks}
%-------------------------------------------------------------------

This appendix specifies the exact task generation procedures for the
three benchmark families, designed to mirror Davis et
al.~(2026)~\cite{davis2026ldt} as closely as possible.

\subsection{Task 1: Propositional SAT}

Random 3-SAT instances generated using Dimitriou--Spirakis clause
sampling. Training: clause-to-variable ratios $\alpha \in [3.0, 4.0]$,
10--20 variables. Test (in-distribution): same ratios and sizes. Test
(OOD): 20--40 variables. Tokenisation: each clause is a token; variables
are embedded as learnable vectors.

\subsection{Task 2: Syllogistic deduction}

Chains $A_1 \subset A_2 \subset \cdots \subset A_K$ encoded as
``All $A_i$ are $A_{i+1}$'' statements. Label: whether
$A_1 \subset A_K$ (always true). Training chain length
$K \in \{2,3,4\}$; test (OOD) $K \in \{5,6,7,8\}$.

\subsection{Task 3: Multi-step numeric reasoning}

Iterative modular addition chains
$x_0 \xrightarrow{+a_1} x_1 \xrightarrow{+a_2} \cdots \xrightarrow{+a_K} x_K \pmod{p}$.
Given $(x_0, a_1, \dots, a_K)$, predict $x_K$. Training: $K \leq 4$,
$p = 97$; test (OOD): $K \in \{5,\dots,10\}$.

\end{document}
